\scanpage{187}
des $G_{\mathbb R}$-modules de dim finie sur $\mathbb R$. Le cône
qui exprime cette effectivité est \uline{engendré}
par les \deleted{$v(s')$} rayons définis par les $v(s')$.
Cela ne fait qu'exprimer \inserted{sous forme ``analytique''} le fait que les
frobenius $f(s')\in G_{\mathbb R}^{\unsure{\xi}(s')}(\mathbb R)$ sont tels
que leurs valeurs propres, dans tte représentation
de $G_{\mathbb R}\deleted{(\unsure{\mathbb R})}^{\unsure{\xi}(s')}$ --- et en fait $\unsure{\overline{\mathrm m}|\simeq}$ de
$G_{\deleted{\mathbb R}\inserted{\mathbb C}}^{\unsure{\xi}(s')}$ --- sont des \uline{entiers} algébriques, ce qui est une
condition arithmétique assez draconienne !
\margin{186}
